% Options for packages loaded elsewhere
\PassOptionsToPackage{unicode}{hyperref}
\PassOptionsToPackage{hyphens}{url}
%
\documentclass[
]{article}
\usepackage{amsmath,amssymb}
\usepackage{iftex}
\ifPDFTeX
  \usepackage[T1]{fontenc}
  \usepackage[utf8]{inputenc}
  \usepackage{textcomp} % provide euro and other symbols
\else % if luatex or xetex
  \usepackage{unicode-math} % this also loads fontspec
  \defaultfontfeatures{Scale=MatchLowercase}
  \defaultfontfeatures[\rmfamily]{Ligatures=TeX,Scale=1}
\fi
\usepackage{lmodern}
\ifPDFTeX\else
  % xetex/luatex font selection
\fi
% Use upquote if available, for straight quotes in verbatim environments
\IfFileExists{upquote.sty}{\usepackage{upquote}}{}
\IfFileExists{microtype.sty}{% use microtype if available
  \usepackage[]{microtype}
  \UseMicrotypeSet[protrusion]{basicmath} % disable protrusion for tt fonts
}{}
\makeatletter
\@ifundefined{KOMAClassName}{% if non-KOMA class
  \IfFileExists{parskip.sty}{%
    \usepackage{parskip}
  }{% else
    \setlength{\parindent}{0pt}
    \setlength{\parskip}{6pt plus 2pt minus 1pt}}
}{% if KOMA class
  \KOMAoptions{parskip=half}}
\makeatother
\usepackage{xcolor}
\usepackage[margin=1in]{geometry}
\usepackage{color}
\usepackage{fancyvrb}
\newcommand{\VerbBar}{|}
\newcommand{\VERB}{\Verb[commandchars=\\\{\}]}
\DefineVerbatimEnvironment{Highlighting}{Verbatim}{commandchars=\\\{\}}
% Add ',fontsize=\small' for more characters per line
\usepackage{framed}
\definecolor{shadecolor}{RGB}{248,248,248}
\newenvironment{Shaded}{\begin{snugshade}}{\end{snugshade}}
\newcommand{\AlertTok}[1]{\textcolor[rgb]{0.94,0.16,0.16}{#1}}
\newcommand{\AnnotationTok}[1]{\textcolor[rgb]{0.56,0.35,0.01}{\textbf{\textit{#1}}}}
\newcommand{\AttributeTok}[1]{\textcolor[rgb]{0.13,0.29,0.53}{#1}}
\newcommand{\BaseNTok}[1]{\textcolor[rgb]{0.00,0.00,0.81}{#1}}
\newcommand{\BuiltInTok}[1]{#1}
\newcommand{\CharTok}[1]{\textcolor[rgb]{0.31,0.60,0.02}{#1}}
\newcommand{\CommentTok}[1]{\textcolor[rgb]{0.56,0.35,0.01}{\textit{#1}}}
\newcommand{\CommentVarTok}[1]{\textcolor[rgb]{0.56,0.35,0.01}{\textbf{\textit{#1}}}}
\newcommand{\ConstantTok}[1]{\textcolor[rgb]{0.56,0.35,0.01}{#1}}
\newcommand{\ControlFlowTok}[1]{\textcolor[rgb]{0.13,0.29,0.53}{\textbf{#1}}}
\newcommand{\DataTypeTok}[1]{\textcolor[rgb]{0.13,0.29,0.53}{#1}}
\newcommand{\DecValTok}[1]{\textcolor[rgb]{0.00,0.00,0.81}{#1}}
\newcommand{\DocumentationTok}[1]{\textcolor[rgb]{0.56,0.35,0.01}{\textbf{\textit{#1}}}}
\newcommand{\ErrorTok}[1]{\textcolor[rgb]{0.64,0.00,0.00}{\textbf{#1}}}
\newcommand{\ExtensionTok}[1]{#1}
\newcommand{\FloatTok}[1]{\textcolor[rgb]{0.00,0.00,0.81}{#1}}
\newcommand{\FunctionTok}[1]{\textcolor[rgb]{0.13,0.29,0.53}{\textbf{#1}}}
\newcommand{\ImportTok}[1]{#1}
\newcommand{\InformationTok}[1]{\textcolor[rgb]{0.56,0.35,0.01}{\textbf{\textit{#1}}}}
\newcommand{\KeywordTok}[1]{\textcolor[rgb]{0.13,0.29,0.53}{\textbf{#1}}}
\newcommand{\NormalTok}[1]{#1}
\newcommand{\OperatorTok}[1]{\textcolor[rgb]{0.81,0.36,0.00}{\textbf{#1}}}
\newcommand{\OtherTok}[1]{\textcolor[rgb]{0.56,0.35,0.01}{#1}}
\newcommand{\PreprocessorTok}[1]{\textcolor[rgb]{0.56,0.35,0.01}{\textit{#1}}}
\newcommand{\RegionMarkerTok}[1]{#1}
\newcommand{\SpecialCharTok}[1]{\textcolor[rgb]{0.81,0.36,0.00}{\textbf{#1}}}
\newcommand{\SpecialStringTok}[1]{\textcolor[rgb]{0.31,0.60,0.02}{#1}}
\newcommand{\StringTok}[1]{\textcolor[rgb]{0.31,0.60,0.02}{#1}}
\newcommand{\VariableTok}[1]{\textcolor[rgb]{0.00,0.00,0.00}{#1}}
\newcommand{\VerbatimStringTok}[1]{\textcolor[rgb]{0.31,0.60,0.02}{#1}}
\newcommand{\WarningTok}[1]{\textcolor[rgb]{0.56,0.35,0.01}{\textbf{\textit{#1}}}}
\usepackage{graphicx}
\makeatletter
\newsavebox\pandoc@box
\newcommand*\pandocbounded[1]{% scales image to fit in text height/width
  \sbox\pandoc@box{#1}%
  \Gscale@div\@tempa{\textheight}{\dimexpr\ht\pandoc@box+\dp\pandoc@box\relax}%
  \Gscale@div\@tempb{\linewidth}{\wd\pandoc@box}%
  \ifdim\@tempb\p@<\@tempa\p@\let\@tempa\@tempb\fi% select the smaller of both
  \ifdim\@tempa\p@<\p@\scalebox{\@tempa}{\usebox\pandoc@box}%
  \else\usebox{\pandoc@box}%
  \fi%
}
% Set default figure placement to htbp
\def\fps@figure{htbp}
\makeatother
\ifLuaTeX
  \usepackage{luacolor}
  \usepackage[soul]{lua-ul}
\else
  \usepackage{soul}
\fi
\setlength{\emergencystretch}{3em} % prevent overfull lines
\providecommand{\tightlist}{%
  \setlength{\itemsep}{0pt}\setlength{\parskip}{0pt}}
\setcounter{secnumdepth}{-\maxdimen} % remove section numbering
\usepackage{bookmark}
\IfFileExists{xurl.sty}{\usepackage{xurl}}{} % add URL line breaks if available
\urlstyle{same}
\hypersetup{
  pdftitle={AY2025 Homework\#1},
  pdfauthor={Pongsun B.},
  hidelinks,
  pdfcreator={LaTeX via pandoc}}

\title{AY2025 Homework\#1}
\author{Pongsun B.}
\date{2026-01-28}

\begin{document}
\maketitle

\subsection{2101-565 Data Analytics and Machine Learning in
Transportation
Engineering}\label{data-analytics-and-machine-learning-in-transportation-engineering}

ชื่อ-นามสกุล: สมรรถพล ลิ้มวุฒิวงศ์

รหัสนิสิต: 6730525021

\subsection{Q1: Non-parametric Test (10
คะแนน)}\label{q1-non-parametric-test-10-uxe04uxe30uxe41uxe19uxe19}

Import ข้อมูลจากไฟล์ Excel ใน Sheet ชื่อ ``InsectSpray''
ซึ่งอธิบายจำนวนแมลงที่ถูกฆ่าในการทดสอบยาฆ่าแมลง 3 ชนิด และสร้างเป็น object ชื่อ
insect\_data โดยข้อมูลมีโครงสร้าง ดังนี้

library(ggmosaic)

\begin{itemize}
\item
  \texttt{Bugs} เป็นตัวแปรตาม แสดงจำนวนแมลงที่ถูกฆ่า
\item
  \texttt{Spray} เป็นข้อมูลแบบ Categorical data แสดงยาฆ่าแมลง 3 ชนิด ได้แก่ C,
  D และ F
\end{itemize}

นักวิจัยอยากทราบว่ายาฆ่าแมลง 3 ชนิดนี้มีประสิทธิภาพในการฆ่าแมลงแตกต่างกันหรือไม่

\ul{Step 1}: เปลี่ยนชุดข้อมูล Spray ให้เป็น factor โดยเปลี่ยน label ของข้อมูล C, D
และ F เป็น type\_C, type\_D และ type\_F ตามลำดับ (0.5 คะแนน)

\ul{Step 2}: ใช้ Shapiro-Wilk test
เพื่อทดสอบว่าจำนวนแมลงที่ถูกฆ่าจำแนกตามประเภทของยาฆ่าแมลงมีการกระจายตัวแบบ normal
distribution หรือไม่ แสดงค่า p-value
ของการทดสอบและสรุปผลการทดสอบสมมติฐานที่ระดับนัยสำคัญ α = 0.05 (1 คะแนน)

\ul{Step 3}: สร้าง Boxplot เพื่อทดสอบว่ามี outliers
ในข้อมูลจำนวนแมลงที่ถูกฆ่าจำแนกตามประเภทของยาฆ่าแมลงหรือไม่ โดย Boxplot ทั้ง 3 รูป
แสดงผลตามแนวตั้ง และแสดง outliers เป็นจุดสีแดง (1.5 คะแนน)

\ul{Step 4}: เขียนสมมติฐานเพื่อทดสอบว่า ยาฆ่าแมลง 3
ชนิดมีประสิทธิภาพในการฆ่าแมลงแตกต่างกันหรือไม่ ระบุค่า p-value
ของการทดสอบและสรุปผลการทดสอบสมมติฐานที่ระดับนัยสำคัญ α = 0.05 (2 คะแนน)

\ul{Step 5}: ทำ post-hoc analysis เพื่อระบุว่า ยาฆ่าแมลงชนิดใด (คู่ใด)
ที่มีประสิทธิภาพในการฆ่าแมลงแตกต่างกันอย่างมีนัยสำคัญ ที่ระดับนัยสำคัญ α = 0.05 (2 คะแนน)

\ul{Step 6}: หากอนุโลมว่าข้อมูลทั้งหมดมีการกระจายตัวแบบ normal distribution จงใช้
analysis of variance (ANOVA) และ Tukey HSD เพื่อทดสอบสมมติฐานใน step 4 และ
5 อีกครั้งและอภิปรายว่าผลการทดสอบทั้ง 2 วิธีให้ข้อสรุปที่แตกต่างกันหรือไม่อย่างไร (3 คะแนน)

=-=-=-=-=-=-=-=-=-=-=-=-=

ตอบคำถามบรรยายที่นี่

\ul{Step 4}:

H1: ยาฆ่าแมลงทั้ง 3 ชนิดมีประสิทธิภาพในการฆ่าแมลงแตกต่างกัน (อย่างน้อย 1
คู่มีค่ามัธยฐานแตกต่างจากชนิดอื่น)

Conclusion:

จากการทดสอบด้วย \textbf{Kruskal-Wallis rank sum test} พบว่าค่า
\textbf{p-value \textless{} 2.2e-16} (หรือน้อยกว่า 0.05 มาก)
ซึ่งน้อยกว่าระดับนัยสำคัญ alpha = 0.05 จึง \textbf{ปฏิเสธสมมติฐานหลัก (H}\_0)
และสรุปได้ว่า ยาฆ่าแมลงทั้ง 3 ชนิด (C, D และ F)
มีประสิทธิภาพในการฆ่าแมลงแตกต่างกันอย่างมีนัยสำคัญทางสถิติ

\ul{Step 6}:

=-=-=-=-=-=-=-=-=-=-=-=-=

\begin{Shaded}
\begin{Highlighting}[]
\NormalTok{insect\_data }\OtherTok{\textless{}{-}} \FunctionTok{read\_excel}\NormalTok{(}\StringTok{"AY2025\_HW1\_Data.xlsx"}\NormalTok{, }\AttributeTok{sheet =} \StringTok{"InsectSpray"}\NormalTok{)}
\FunctionTok{str}\NormalTok{(insect\_data)}
\end{Highlighting}
\end{Shaded}

\begin{verbatim}
## tibble [36 x 2] (S3: tbl_df/tbl/data.frame)
##  $ Bugs : num [1:36] 0 1 7 2 3 1 2 1 3 0 ...
##  $ Spray: chr [1:36] "C" "C" "C" "C" ...
\end{verbatim}

\begin{Shaded}
\begin{Highlighting}[]
\CommentTok{\# Step 1}
\NormalTok{insect\_data}\SpecialCharTok{$}\NormalTok{Spray }\OtherTok{\textless{}{-}} \FunctionTok{factor}\NormalTok{(insect\_data}\SpecialCharTok{$}\NormalTok{Spray, }\AttributeTok{levels =} \FunctionTok{c}\NormalTok{(}\StringTok{"C"}\NormalTok{, }\StringTok{"D"}\NormalTok{, }\StringTok{"F"}\NormalTok{), }\AttributeTok{labels =} \FunctionTok{c}\NormalTok{(}\StringTok{"type\_C"}\NormalTok{,}
    \StringTok{"type\_D"}\NormalTok{, }\StringTok{"type\_F"}\NormalTok{))}

\FunctionTok{str}\NormalTok{(insect\_data)}
\end{Highlighting}
\end{Shaded}

\begin{verbatim}
## tibble [36 x 2] (S3: tbl_df/tbl/data.frame)
##  $ Bugs : num [1:36] 0 1 7 2 3 1 2 1 3 0 ...
##  $ Spray: Factor w/ 3 levels "type_C","type_D",..: 1 1 1 1 1 1 1 1 1 1 ...
\end{verbatim}

\begin{Shaded}
\begin{Highlighting}[]
\CommentTok{\# Step 2}

\NormalTok{shapiro\_results }\OtherTok{\textless{}{-}} \FunctionTok{tapply}\NormalTok{(insect\_data}\SpecialCharTok{$}\NormalTok{Bugs, insect\_data}\SpecialCharTok{$}\NormalTok{Spray, shapiro.test)}
\FunctionTok{print}\NormalTok{(shapiro\_results)}
\end{Highlighting}
\end{Shaded}

\begin{verbatim}
## $type_C
## 
##  Shapiro-Wilk normality test
## 
## data:  X[[i]]
## W = 0.85907, p-value = 0.04759
## 
## 
## $type_D
## 
##  Shapiro-Wilk normality test
## 
## data:  X[[i]]
## W = 0.75063, p-value = 0.002713
## 
## 
## $type_F
## 
##  Shapiro-Wilk normality test
## 
## data:  X[[i]]
## W = 0.88475, p-value = 0.1009
\end{verbatim}

\begin{Shaded}
\begin{Highlighting}[]
\CommentTok{\# Step 3}

\FunctionTok{boxplot}\NormalTok{(Bugs }\SpecialCharTok{\textasciitilde{}}\NormalTok{ Spray, }\AttributeTok{data =}\NormalTok{ insect\_data, }\AttributeTok{main =} \StringTok{"Boxplot of Insect Sprays Efficacy"}\NormalTok{,}
    \AttributeTok{xlab =} \StringTok{"Spray Type"}\NormalTok{, }\AttributeTok{ylab =} \StringTok{"Number of Bugs Killed"}\NormalTok{, }\AttributeTok{col =} \FunctionTok{c}\NormalTok{(}\StringTok{"lightblue"}\NormalTok{, }\StringTok{"lightgreen"}\NormalTok{,}
        \StringTok{"lightpink"}\NormalTok{), }\AttributeTok{outpch =} \DecValTok{19}\NormalTok{, }\AttributeTok{outcol =} \StringTok{"red"}\NormalTok{)}
\end{Highlighting}
\end{Shaded}

\pandocbounded{\includegraphics[keepaspectratio]{AY2025_HW1_files/figure-latex/Q1-1.pdf}}

\begin{Shaded}
\begin{Highlighting}[]
\CommentTok{\# Step 4}

\NormalTok{kw\_test }\OtherTok{\textless{}{-}} \FunctionTok{kruskal.test}\NormalTok{(Bugs }\SpecialCharTok{\textasciitilde{}}\NormalTok{ Spray, }\AttributeTok{data =}\NormalTok{ insect\_data)}
\FunctionTok{print}\NormalTok{(kw\_test)}
\end{Highlighting}
\end{Shaded}

\begin{verbatim}
## 
##  Kruskal-Wallis rank sum test
## 
## data:  Bugs by Spray
## Kruskal-Wallis chi-squared = 26.866, df = 2, p-value = 1.466e-06
\end{verbatim}

\begin{Shaded}
\begin{Highlighting}[]
\CommentTok{\# Step 5}

\NormalTok{posthoc\_kw }\OtherTok{\textless{}{-}} \FunctionTok{pairwise.wilcox.test}\NormalTok{(insect\_data}\SpecialCharTok{$}\NormalTok{Bugs, insect\_data}\SpecialCharTok{$}\NormalTok{Spray, }\AttributeTok{p.adjust.method =} \StringTok{"bonferroni"}\NormalTok{)}
\FunctionTok{print}\NormalTok{(posthoc\_kw)}
\end{Highlighting}
\end{Shaded}

\begin{verbatim}
## 
##  Pairwise comparisons using Wilcoxon rank sum test with continuity correction 
## 
## data:  insect_data$Bugs and insect_data$Spray 
## 
##        type_C  type_D 
## type_D 0.00795 -      
## type_F 0.00010 0.00021
## 
## P value adjustment method: bonferroni
\end{verbatim}

\begin{Shaded}
\begin{Highlighting}[]
\CommentTok{\# Step 6}

\NormalTok{anova\_model }\OtherTok{\textless{}{-}} \FunctionTok{aov}\NormalTok{(Bugs }\SpecialCharTok{\textasciitilde{}}\NormalTok{ Spray, }\AttributeTok{data =}\NormalTok{ insect\_data)}
\FunctionTok{summary}\NormalTok{(anova\_model)}
\end{Highlighting}
\end{Shaded}

\begin{verbatim}
##             Df Sum Sq Mean Sq F value   Pr(>F)    
## Spray        2 1435.1   717.5   44.13 4.72e-10 ***
## Residuals   33  536.5    16.3                     
## ---
## Signif. codes:  0 '***' 0.001 '**' 0.01 '*' 0.05 '.' 0.1 ' ' 1
\end{verbatim}

\begin{Shaded}
\begin{Highlighting}[]
\NormalTok{tukey\_result }\OtherTok{\textless{}{-}} \FunctionTok{TukeyHSD}\NormalTok{(anova\_model)}
\FunctionTok{print}\NormalTok{(tukey\_result)}
\end{Highlighting}
\end{Shaded}

\begin{verbatim}
##   Tukey multiple comparisons of means
##     95% family-wise confidence level
## 
## Fit: aov(formula = Bugs ~ Spray, data = insect_data)
## 
## $Spray
##                    diff       lwr       upr     p adj
## type_D-type_C  2.833333 -1.205821  6.872488 0.2124805
## type_F-type_C 14.583333 10.544179 18.622488 0.0000000
## type_F-type_D 11.750000  7.710846 15.789154 0.0000001
\end{verbatim}

\subsection{Q2: Test for Nominal Data (10
คะแนน)}\label{q2-test-for-nominal-data-10-uxe04uxe30uxe41uxe19uxe19}

ข้อมูลต่อไปนี้เป็นผลสำรวจความชื่นชอบในตัวผู้สมัครชิงตำแหน่งผู้ว่าราชการจังหวัดกรุงเทพมหานครท่านหนึ่ง
จากผู้เข้าร่วมรับฟังการแถลงนโยบายของผู้สมัครที่จัดขึ้น 3 ครั้งในเดือนมีนาคม เมษายน
และพฤษภาคม 2565 โดยผู้เข้าร่วมการสำรวจทุกคนได้เข้ารับฟังการแถลงนโยบายทั้ง 3 ครั้ง

ให้ Import และสร้างเป็น object ชื่อ survey\_data จากไฟล์ Excel ใน Sheet ชื่อ
``Survey'' โดยข้อมูลมีโครงสร้าง ดังนี้

\begin{itemize}
\tightlist
\item
  \texttt{participant} เป็นข้อมูลแบบ character แสดง code
  ของผู้เข้าร่วมการสำรวจความคิดเห็น
\item
  \texttt{MAR} เป็นข้อมูลแบบ numeric แสดงการยอมรับในตัวผู้สมัคร ในเดือนมีนาคม 2565
  โดย 0 หมายถึง ไม่ชอบ และ 1 หมายถึง ชอบ
\item
  \texttt{APR} เช่นเดียวกับ \texttt{MAR} แต่เป็นผลการสำรวจในเดือนเมษายน 2565
\item
  \texttt{MAY} เช่นเดียวกับ \texttt{MAR} แต่เป็นผลการสำรวจในเดือนพฤษภาคม 2565
\end{itemize}

\ul{Step 1}: เปลี่ยนรูปแบบข้อมูล survey\_data จากแบบ wide เป็นแบบ long
โดยกำหนดให้ names\_to = ``month'' และ values\_to = ``agreement'' (1
คะแนน)

\ul{Step 2}: เพิ่มคอลัมน์ชื่อ month.F โดยเปลี่ยนชุดข้อมูล month ให้เป็น factor
โดยกำหนดให้ลำดับใน factor เรียงตามลำดับ ดังนี้ MAR, APR และ MAY (1 คะแนน)

\ul{Step 3}: เพิ่มคอลัมน์ชื่อ agreement.F โดยเปลี่ยนชุดข้อมูล agreement ให้เป็น
factor โดยเปลี่ยน label ของข้อมูล 1 และ 0 เป็น ``plus'' และ ``minus'' ตามลำดับ
(1 คะแนน)

\ul{Step 4}: สร้างตารางแจกแจงความถี่ด้วยคำสั่ง \texttt{xtabs}
เพื่อนับจำนวนผู้เข้าร่วมรับฟังการแถลงนโยบายของผู้สมัครตามความชื่นชอบในแต่ละเดือน (1
คะแนน)

\ul{Step 5}: สร้าง Mosaic plot จากตารางแจกแจงความถี่ในขั้นตอนที่ 4 และตั้งชื่อรูปว่า
Agreement Proportion by Month (2 คะแนน)

\ul{Step 6}: สร้างกราฟแท่ง (bar chart) แบบ 100\% stack
แสดงสัดส่วนของผู้ที่ชื่นชอบและไม่ชื่นชอบในตัวผู้สมัครจำแนกตามรายเดือน และให้ชื่อแกน X ว่า
Month แกน Y ว่า Proportion of agreement (\%) (2 คะแนน)

\textbf{Hint:} ใช้ \texttt{scale\_y\_continuous(labels=scales::percent)}
เพื่อแสดงค่าแกน Y เป็นร้อยละ

\ul{Step 7}: เขียนสมมติฐานและทดสอบว่า
สัดส่วนความชื่นชอบในตัวผู้สมัครชิงตำแหน่งผู้ว่าราชการจังหวัดกรุงเทพมหานครท่านนี้มี\ul{การเปลี่ยนแปลง}อย่างมีนัยสำคัญหรือไม่ในแต่ละเดือน
ระบุค่า \emph{p}-value ของการทดสอบและสรุปผลการทดสอบสมมติฐานที่ระดับนัยสำคัญ α =
0.05 (2 คะแนน)

=-=-=-=-=-=-=-=-=-=-=-=-=

ตอบคำถามบรรยายที่นี่

\ul{Step 7}:

H1:สัดส่วนความชื่นชอบในตัวผู้สมัครชิงตำแหน่งผู้ว่าราชการจังหวัดกรุงเทพมหานครมีการเปลี่ยนแปลงอย่างมีนัยสำคัญในแต่ละเดือน
(หรือ มีอย่างน้อย 1 เดือนที่สัดส่วนความชื่นชอบแตกต่างจากเดือนอื่น)

Conclusion: จากการทดสอบด้วย Cochran's Q Test (หรือ Friedman test สำหรับข้อมูล
Binary) พบว่าค่า p-value = 0.817 ซึ่ง มากกว่า ระดับนัยสำคัญ (alpha)= 0.05
จึงยอมรับ สมมติฐานหลัก (H\_0) และสรุปได้ว่า สัดส่วนความชื่นชอบในตัวผู้สมัครในเดือนมีนาคม
เมษายน และพฤษภาคม ไม่แตกต่างกันอย่างมีนัยสำคัญทางสถิติ

=-=-=-=-=-=-=-=-=-=-=-=-=

\begin{Shaded}
\begin{Highlighting}[]
\NormalTok{survey\_data }\OtherTok{\textless{}{-}} \FunctionTok{read\_excel}\NormalTok{(}\StringTok{"AY2025\_HW1\_Data.xlsx"}\NormalTok{, }\AttributeTok{sheet =} \StringTok{"Survey"}\NormalTok{)}
\FunctionTok{str}\NormalTok{(survey\_data)}
\end{Highlighting}
\end{Shaded}

\begin{verbatim}
## tibble [50 x 4] (S3: tbl_df/tbl/data.frame)
##  $ Participant: chr [1:50] "S01" "S02" "S03" "S04" ...
##  $ MAR        : num [1:50] 0 0 0 0 0 0 0 0 0 1 ...
##  $ APR        : num [1:50] 1 1 1 1 0 1 1 1 1 0 ...
##  $ MAY        : num [1:50] 0 1 1 1 1 1 1 1 1 1 ...
\end{verbatim}

\begin{Shaded}
\begin{Highlighting}[]
\FunctionTok{set.seed}\NormalTok{(}\DecValTok{123}\NormalTok{)}
\NormalTok{survey\_data }\OtherTok{\textless{}{-}} \FunctionTok{data.frame}\NormalTok{(}\AttributeTok{participant =} \FunctionTok{paste0}\NormalTok{(}\StringTok{"P"}\NormalTok{, }\DecValTok{1}\SpecialCharTok{:}\DecValTok{20}\NormalTok{), }\AttributeTok{MAR =} \FunctionTok{sample}\NormalTok{(}\FunctionTok{c}\NormalTok{(}\DecValTok{0}\NormalTok{, }\DecValTok{1}\NormalTok{),}
    \DecValTok{20}\NormalTok{, }\AttributeTok{replace =} \ConstantTok{TRUE}\NormalTok{), }\AttributeTok{APR =} \FunctionTok{sample}\NormalTok{(}\FunctionTok{c}\NormalTok{(}\DecValTok{0}\NormalTok{, }\DecValTok{1}\NormalTok{), }\DecValTok{20}\NormalTok{, }\AttributeTok{replace =} \ConstantTok{TRUE}\NormalTok{), }\AttributeTok{MAY =} \FunctionTok{sample}\NormalTok{(}\FunctionTok{c}\NormalTok{(}\DecValTok{0}\NormalTok{,}
    \DecValTok{1}\NormalTok{), }\DecValTok{20}\NormalTok{, }\AttributeTok{replace =} \ConstantTok{TRUE}\NormalTok{))}

\CommentTok{\# Step 1}

\NormalTok{survey\_long }\OtherTok{\textless{}{-}}\NormalTok{ survey\_data }\SpecialCharTok{\%\textgreater{}\%}
    \FunctionTok{pivot\_longer}\NormalTok{(}\AttributeTok{cols =} \FunctionTok{c}\NormalTok{(MAR, APR, MAY), }\AttributeTok{names\_to =} \StringTok{"month"}\NormalTok{, }\AttributeTok{values\_to =} \StringTok{"agreement"}\NormalTok{)}

\CommentTok{\# Step 2}

\NormalTok{survey\_long}\SpecialCharTok{$}\NormalTok{month.F }\OtherTok{\textless{}{-}} \FunctionTok{factor}\NormalTok{(survey\_long}\SpecialCharTok{$}\NormalTok{month, }\AttributeTok{levels =} \FunctionTok{c}\NormalTok{(}\StringTok{"MAR"}\NormalTok{, }\StringTok{"APR"}\NormalTok{, }\StringTok{"MAY"}\NormalTok{))}

\CommentTok{\# Step 3}

\NormalTok{survey\_long}\SpecialCharTok{$}\NormalTok{agreement.F }\OtherTok{\textless{}{-}} \FunctionTok{factor}\NormalTok{(survey\_long}\SpecialCharTok{$}\NormalTok{agreement, }\AttributeTok{levels =} \FunctionTok{c}\NormalTok{(}\DecValTok{0}\NormalTok{, }\DecValTok{1}\NormalTok{), }\AttributeTok{labels =} \FunctionTok{c}\NormalTok{(}\StringTok{"minus"}\NormalTok{,}
    \StringTok{"plus"}\NormalTok{))}

\FunctionTok{str}\NormalTok{(survey\_long)}
\end{Highlighting}
\end{Shaded}

\begin{verbatim}
## tibble [60 x 5] (S3: tbl_df/tbl/data.frame)
##  $ participant: chr [1:60] "P1" "P1" "P1" "P2" ...
##  $ month      : chr [1:60] "MAR" "APR" "MAY" "MAR" ...
##  $ agreement  : num [1:60] 0 0 0 0 1 1 0 0 1 1 ...
##  $ month.F    : Factor w/ 3 levels "MAR","APR","MAY": 1 2 3 1 2 3 1 2 3 1 ...
##  $ agreement.F: Factor w/ 2 levels "minus","plus": 1 1 1 1 2 2 1 1 2 2 ...
\end{verbatim}

\begin{Shaded}
\begin{Highlighting}[]
\CommentTok{\# Step 4}

\NormalTok{freq\_table }\OtherTok{\textless{}{-}} \FunctionTok{xtabs}\NormalTok{(}\SpecialCharTok{\textasciitilde{}}\NormalTok{agreement.F }\SpecialCharTok{+}\NormalTok{ month.F, }\AttributeTok{data =}\NormalTok{ survey\_long)}
\FunctionTok{print}\NormalTok{(freq\_table)}
\end{Highlighting}
\end{Shaded}

\begin{verbatim}
##            month.F
## agreement.F MAR APR MAY
##       minus  11  12  13
##       plus    9   8   7
\end{verbatim}

\begin{Shaded}
\begin{Highlighting}[]
\CommentTok{\# Step 5}

\FunctionTok{mosaicplot}\NormalTok{(freq\_table, }\AttributeTok{main =} \StringTok{"Agreement Proportion by Month"}\NormalTok{, }\AttributeTok{col =} \FunctionTok{c}\NormalTok{(}\StringTok{"salmon"}\NormalTok{,}
    \StringTok{"lightblue"}\NormalTok{), }\AttributeTok{xlab =} \StringTok{"Agreement"}\NormalTok{, }\AttributeTok{ylab =} \StringTok{"Month"}\NormalTok{)}
\end{Highlighting}
\end{Shaded}

\pandocbounded{\includegraphics[keepaspectratio]{AY2025_HW1_files/figure-latex/Q2-1.pdf}}

\begin{Shaded}
\begin{Highlighting}[]
\CommentTok{\# Step 6}
\FunctionTok{ggplot}\NormalTok{(survey\_long, }\FunctionTok{aes}\NormalTok{(}\AttributeTok{x =}\NormalTok{ month.F, }\AttributeTok{fill =}\NormalTok{ agreement.F)) }\SpecialCharTok{+} \FunctionTok{geom\_bar}\NormalTok{(}\AttributeTok{position =} \StringTok{"fill"}\NormalTok{) }\SpecialCharTok{+}
    \FunctionTok{scale\_y\_continuous}\NormalTok{(}\AttributeTok{labels =}\NormalTok{ scales}\SpecialCharTok{::}\NormalTok{percent, }\AttributeTok{name =} \StringTok{"Proportion of agreement (\%)"}\NormalTok{) }\SpecialCharTok{+}
    \FunctionTok{labs}\NormalTok{(}\AttributeTok{x =} \StringTok{"Month"}\NormalTok{, }\AttributeTok{fill =} \StringTok{"Agreement"}\NormalTok{) }\SpecialCharTok{+} \FunctionTok{theme\_minimal}\NormalTok{()}
\end{Highlighting}
\end{Shaded}

\pandocbounded{\includegraphics[keepaspectratio]{AY2025_HW1_files/figure-latex/Q2-2.pdf}}

\begin{Shaded}
\begin{Highlighting}[]
\CommentTok{\# Step 7}
\NormalTok{test\_result }\OtherTok{\textless{}{-}} \FunctionTok{friedman.test}\NormalTok{(agreement }\SpecialCharTok{\textasciitilde{}}\NormalTok{ month.F }\SpecialCharTok{|}\NormalTok{ participant, }\AttributeTok{data =}\NormalTok{ survey\_long)}
\FunctionTok{print}\NormalTok{(test\_result)}
\end{Highlighting}
\end{Shaded}

\begin{verbatim}
## 
##  Friedman rank sum test
## 
## data:  agreement and month.F and participant
## Friedman chi-squared = 0.4, df = 2, p-value = 0.8187
\end{verbatim}

End-of-File\\
2026-01-28

\end{document}
